\section{Algorithm}
In this section we denote the target state as $\ket{\psi_{target}}$ 
and the state to be verified as $\ket{\psi_{lab}}$. 
Let's directly present the algorihm first:

\begin{algorithm}[H]
    \vspace{0.3em}
    \textbf{Input:} Full classical description of $\ket{\psi_{target}}$ 
    and one copy of an unknown state $\ket{\psi_{lab}}$. \\
    \textbf{Output:} \textsc{Accept} or \textsc{Reject}.
    \begin{algorithmic}[1]
    \State \textbf{Sample} $k\in[n]$ uniformly at random.
    \State \textbf{Measure} the first $k-1$ qubits of $\ket{\psi_{lab}}$ 
    in the computational basis, obtaining~$x\in\{0,1\}^{k-1}$. 
    Let the reduced state of the last $n-k+1$ qubits be $\ket{\psi_{lab}^x}$ and 
    that of $\ket{\psi_{target}}$ be $\ket{\psi_{target}^x}$ when the same outcome is obtained.
    \State Let $\ket{\psi_{target}^x} = \alpha\ket{0}\ket{\psi_{target}^{x0}} + \beta\ket{1}\ket{\psi_{target}^{x1}}$, 
    where $\ket{\psi_{target}^{x0}}$ and $\ket{\psi_{target}^{x1}}$ are states on the last $n-k$ qubits.
    
    \noindent Define $\mathcal{T}_x$ as the basis where both $\ket{\psi_{target}^{x0}}$ and $\ket{\psi_{target}^{x1}}$ are both 
    phase states (*). 
    \textbf{Measure} the last $n-k$ qubits of $\ket{\psi_{lab}^x}$ 
    in the basis $\mathcal{T}_x$, obtaining~$\ell$.
    \State Let $\ket{\psi_{lab}'}$ be the resulting $1$-qubit state; 
    let $\ket{\psi_{target}'}$ denote $\ket{\psi_{target}}$ conditioned on outcomes $x, \ell$.
    \noindent \textbf{Measure} $\ket{\psi_{lab}'}$ in a basis containing $\ket{\psi_{target}'}$,
    and \textbf{Output} \textsc{Accept} iff the outcome is $\ket{\psi_{target}'}$.
    \end{algorithmic}
    \caption{\textsc{Certify}$(\ket{\psi_{lab}}, \ket{\psi_{target}})$}
\end{algorithm}

(*)Note: A state $\ket{\phi}$ on $m$ qubits is called a \emph{phase state} if it can be written as
\[
    \ket{\phi} = \frac{1}{\sqrt{2^m}} \sum_{y \in \{0,1\}^m} e^{i\theta_y} \ket{y},
\]
for some phases $\theta_y \in \mathbb{R}$, i.e.\ every computational basis amplitude has the same magnitude $1/\sqrt{2^m}$.
The existence of the basis $\mathcal{T}_x$ need to be proven, and we will do so in the next subsection. 

It suffices to show that the acceptance and rejection probabilities of the above algorithm are as claimed in section 1.

\subsection{Acceptance and rejection probabilities}

We refer to Steps 3--4 of our algorithm as the \emph{subtest} 
performed on the $(n-k+1)$-qubit states $\ket{\psi_{target}^x}$ and 
$\ket{\psi_{lab}^x}$ with the basis $\mathcal{T}_x$, and write
\[
    \textsc{SubTest}_{\mathcal{T}_x}\!\left(\ket{\psi_{lab}^x}:\ket{\psi_{target}^x}\right)
\]
for the corresponding two-outcome experiment.
The key to analyzing our algorithm is to compare the subtest acceptance 
probability to the quantity defined below.

\begin{definition}[Fidelity gap]
    For $(n-k+1)$-qubit states $\ket{\psi_{lab}^x}$ and $\ket{\psi_{target}^x}$, 
    we define their \emph{fidelity gap} to be
    \[
        \Delta(\ket{\psi_{lab}^x}:\ket{\psi_{target}^x}) \coloneqq \mathop{\mathbb{E}}_{b}\!\left[\,\abs{\braket{\psi_{target}^{xb}}{\psi_{lab}^{xb}}}^2\,\right] - \abs{\braket{\psi_{target}^x}{\psi_{lab}^x}}^2.
    \]
    Here the expectation is over the random outcome $b \in \{0,1\}$ 
    obtained when measuring the first qubit of $\ket{\psi_{lab}^x}$ in the computational basis.
\end{definition}

\begin{proposition}\label{rem:nonneg}
    $\Delta(\ket{\psi_{lab}^x}:\ket{\psi_{target}^x}) \geq 0$ always holds. 
\end{proposition}

\begin{proof}
    Decompose $\ket{\psi_{target}^x}$ and $\ket{\psi_{lab}^x}$ according to their first qubit in the computational basis:
\[
    \ket{\psi_{target}^x} = \alpha\,\ket{0}\otimes\ket{\psi_{target}^{x0}} + \beta\,\ket{1}\otimes\ket{\psi_{target}^{x1}}, \qquad
    \ket{\psi_{lab}^x} = \gamma\,\ket{0}\otimes\ket{\psi_{lab}^{x0}} + \delta\,\ket{1}\otimes\ket{\psi_{lab}^{x1}},
\]
where $\alpha, \beta \in \mathbb{C}$ are as in 
the algorithm statement and $\gamma, \delta \in \mathbb{C}$ 
are the analogous coefficients for $\ket{\psi_{lab}^x}$.
Taking the inner product of the two decompositions gives
\[
    \braket{\psi_{target}^x}{\psi_{lab}^x}
    = \overline{\alpha}\gamma\,\braket{\psi_{target}^{x0}}{\psi_{lab}^{x0}}
    + \overline{\beta}\delta\,\braket{\psi_{target}^{x1}}{\psi_{lab}^{x1}}.
\]
Applying Cauchy--Schwarz $\abs{c_0 d_0 + c_1 d_1}^2 \leq (\abs{c_0}^2 + \abs{c_1}^2)(\abs{d_0}^2 + \abs{d_1}^2)$ with $c_0 = \overline{\alpha}$, $c_1 = \overline{\beta}$, $d_0 = \gamma\,\braket{\psi_{target}^{x0}}{\psi_{lab}^{x0}}$, and $d_1 = \delta\,\braket{\psi_{target}^{x1}}{\psi_{lab}^{x1}}$ yields
\[
    \abs{\braket{\psi_{target}^x}{\psi_{lab}^x}}^2
    \leq (\abs{\alpha}^2 + \abs{\beta}^2)\left(\abs{\gamma}^2\,\abs{\braket{\psi_{target}^{x0}}{\psi_{lab}^{x0}}}^2 + \abs{\delta}^2\,\abs{\braket{\psi_{target}^{x1}}{\psi_{lab}^{x1}}}^2\right).
\]
The first factor equals $1$, and the second factor is exactly $\mathop{\mathbb{E}}_b\!\left[\abs{\braket{\psi_{target}^{xb}}{\psi_{lab}^{xb}}}^2\right]$ since $\Pr[b=0] = \abs{\gamma}^2$ and $\Pr[b=1] = \abs{\delta}^2$. Hence
\[
    \abs{\braket{\psi_{target}^x}{\psi_{lab}^x}}^2 \leq \mathop{\mathbb{E}}_b\!\left[\abs{\braket{\psi_{target}^{xb}}{\psi_{lab}^{xb}}}^2\right],
\]
which is precisely $\Delta(\ket{\psi_{lab}^x}:\ket{\psi_{target}^x}) \geq 0$.
\end{proof}

\begin{theorem}\label{thm:subtest}
    With $\mathcal{T}_x$ being a DT basis in which $\ket{\psi_{target}^{x0}}$ and $\ket{\psi_{target}^{x1}}$ are both phase states, the following holds:
    \begin{align}
        \Pr\!\left[\textsc{SubTest}_{\mathcal{T}_x}\!\left(\ket{\psi_{lab}^x}:\ket{\psi_{target}^x}\right) \text{ rejects}\right]
            &\geq \Delta(\ket{\psi_{lab}^x}:\ket{\psi_{target}^x}), \label{eqn:reject}\\
        \Pr\!\left[\textsc{SubTest}_{\mathcal{T}_x}\!\left(\ket{\psi_{lab}^x}:\ket{\psi_{target}^x}\right) \text{ accepts}\right]
            &\geq \abs{\braket{\psi_{target}^x}{\psi_{lab}^x}}^2. \label{eqn:accept}
    \end{align}
\end{theorem}

\begin{proof}[Proof of Theorem~\ref{thm:subtest}]
We reuse the decomposition from the proof of Proposition~\ref{rem:nonneg}:
\[
    \ket{\psi_{target}^x} = \alpha\,\ket{0}\otimes\ket{\psi_{target}^{x0}} + \beta\,\ket{1}\otimes\ket{\psi_{target}^{x1}}, \qquad
    \ket{\psi_{lab}^x} = \gamma\,\ket{0}\otimes\ket{\psi_{lab}^{x0}} + \delta\,\ket{1}\otimes\ket{\psi_{lab}^{x1}}.
\]
Since both $\ket{\psi_{target}^{x0}}$ and $\ket{\psi_{target}^{x1}}$ are phase states in $\mathcal{T}_x$, there is a unitary $D$ diagonal in $\mathcal{T}_x$ such that $\ket{\psi_{target}^{x1}} = D\ket{\psi_{target}^{x0}}$. We write the basis vectors of $\mathcal{T}_x$ as $\ket{\ell}$ for $\ell \in [2^{n-k}]$, 
and let $\zeta_\ell$ be the corresponding diagonal entry of $D$ . Let
\[
    A \coloneqq \braket{\psi_{target}^{x0}}{\psi_{lab}^{x0}}, \qquad B \coloneqq \braket{\psi_{target}^{x1}}{\psi_{lab}^{x1}}.
\]

\medskip
\noindent \textbf{Claim:}
\begin{align}
    \Pr\!\left[\textsc{SubTest}_{\mathcal{T}_x}\!\left(\ket{\psi_{lab}^x}:\ket{\psi_{target}^x}\right) \text{ accepts}\right]
        &= \norm{\overline{\alpha}\gamma\,\ket{\psi_{lab}^{x0}} + \overline{\beta}\delta\,D^\dagger\ket{\psi_{lab}^{x1}}}^2, \label{eqn:acc-formula}\\
    \Pr\!\left[\textsc{SubTest}_{\mathcal{T}_x}\!\left(\ket{\psi_{lab}^x}:\ket{\psi_{target}^x}\right) \text{ rejects}\right]
        &= \norm{\beta\gamma\,\ket{\psi_{lab}^{x0}} - \alpha\delta\,D^\dagger\ket{\psi_{lab}^{x1}}}^2. \label{eqn:rej-formula}
\end{align}

\begin{proof}[Proof of claim]
Let $f^0_\ell \coloneqq \braket{\ell}{\psi_{lab}^{x0}}$, $f^1_\ell \coloneqq \braket{\ell}{\psi_{lab}^{x1}}$, and $c_\ell \coloneqq \braket{\ell}{\psi_{target}^{x0}}$. The phase-state property gives $\abs{c_\ell}^2 = 1/2^{n-k}$ and $\braket{\ell}{\psi_{target}^{x1}} = \zeta_\ell c_\ell$. Measuring the last $n-k$ qubits of $\ket{\psi_{lab}^x}$ in basis $\mathcal{T}_x$ yields outcome $\ell$ with probability $p_\ell = \abs{\gamma f^0_\ell}^2 + \abs{\delta f^1_\ell}^2$. 
After that, the reduced 1-qubit state is: 

\[
\ket{\psi_{lab}'} = (\gamma f^0_\ell\ket{0} + \delta f^1_\ell\ket{1})/\sqrt{p_\ell}
\]
The analogous calculation on $\ket{\psi_{target}^x}$ 
shows that up to a global phase $\ket{\psi_{target}'} = \alpha\ket{0} + \beta\zeta_\ell\ket{1}$.
Hence the conditional acceptance probability given outcome $\ell$ is
\[
    \abs{\braket{\psi_{target}'}{\psi_{lab}'}}^2 = \frac{\abs{\overline{\alpha}\gamma f^0_\ell + \overline{\beta}\,\overline{\zeta_\ell}\,\delta f^1_\ell}^2}{p_\ell}
\]
Multiplying by $p_\ell$, summing over $\ell$, 
and using $\overline{\zeta_\ell}\,f^1_\ell = \braket{\ell}{D^\dagger\psi_{lab}^{x1}}$ gives
\[
    \Pr[\textsc{SubTest}\text{ accepts}] = \sum_\ell \abs{\overline{\alpha}\gamma f^0_\ell + \overline{\beta}\,\overline{\zeta_\ell}\,\delta f^1_\ell}^2 = \norm{\overline{\alpha}\gamma\,\ket{\psi_{lab}^{x0}} + \overline{\beta}\delta\,D^\dagger\ket{\psi_{lab}^{x1}}}^2,
\]
establishing \eqref{eqn:acc-formula}.

For \eqref{eqn:rej-formula}, since $\Pr[\text{accepts}] + \Pr[\text{rejects}] = 1$ it suffices to verify that the two right-hand sides of \eqref{eqn:acc-formula} and \eqref{eqn:rej-formula} sum to $1$. Expanding,
\[
    \norm{\overline{\alpha}\gamma\,\ket{\psi_{lab}^{x0}} + \overline{\beta}\delta\,D^\dagger\ket{\psi_{lab}^{x1}}}^2 + \norm{\beta\gamma\,\ket{\psi_{lab}^{x0}} - \alpha\delta\,D^\dagger\ket{\psi_{lab}^{x1}}}^2 = (\abs{\alpha}^2 + \abs{\beta}^2)(\abs{\gamma}^2 + \abs{\delta}^2) = 1,
\]
since the cross terms cancel and $\norm{\ket{\psi_{lab}^{x0}}} = \norm{D^\dagger\ket{\psi_{lab}^{x1}}} = 1$.
\end{proof}

It remains to show $\abs{\braket{\psi_{target}^x}{\psi_{lab}^x}}^2 \leq \norm{\overline{\alpha}\gamma\,\ket{\psi_{lab}^{x0}} + \overline{\beta}\delta\,D^\dagger\ket{\psi_{lab}^{x1}}}^2$ 
and $\Delta(\ket{\psi_{lab}^x}:\ket{\psi_{target}^x}) \leq \norm{\beta\gamma\,\ket{\psi_{lab}^{x0}} - \alpha\delta\,D^\dagger\ket{\psi_{lab}^{x1}}}^2$.
\medskip

Since $\ket{\psi_{target}^{x1}} = D\ket{\psi_{target}^{x0}}$, we have $\bra{\psi_{target}^{x1}} = \bra{\psi_{target}^{x0}}D^\dagger$, and so
\[
    \braket{\psi_{target}^x}{\psi_{lab}^x} = \overline{\alpha}\gamma A + \overline{\beta}\delta B = \bra{\psi_{target}^{x0}}\!\left(\overline{\alpha}\gamma\,\ket{\psi_{lab}^{x0}} + \overline{\beta}\delta\,D^\dagger\ket{\psi_{lab}^{x1}}\right).
\]
The Cauchy--Schwarz inequality $\abs{\braket{u}{w}}^2 \leq \norm{w}^2$ for unit vector $\ket{u}$ ($\ket{u} = \ket{\psi^{x0}_{target}}$ and $\ket{w} = \overline{\alpha}\gamma\,\ket{\psi_{lab}^{x0}} + \overline{\beta}\delta\,D^\dagger\ket{\psi_{lab}^{x1}}$) then gives
\[
    \abs{\braket{\psi_{target}^x}{\psi_{lab}^x}}^2 \leq \norm{\overline{\alpha}\gamma\,\ket{\psi_{lab}^{x0}} + \overline{\beta}\delta\,D^\dagger\ket{\psi_{lab}^{x1}}}^2 \stackrel{\eqref{eqn:acc-formula}}{=} \Pr[\textsc{SubTest}\text{ accepts}],
\]
which is \eqref{eqn:accept}.

\medskip
On the other hand, we first re-express the fidelity gap as
\begin{equation}\label{eqn:delta-reexp}
    \Delta(\ket{\psi_{lab}^x}:\ket{\psi_{target}^x}) = \abs{\beta\gamma A - \alpha\delta B}^2.
\end{equation}
Indeed, since $\Pr[b=0] = \abs{\gamma}^2$ and $\Pr[b=1] = \abs{\delta}^2$,
\begin{align*}
    \Delta(\ket{\psi_{lab}^x}:\ket{\psi_{target}^x})
        &= \abs{\gamma}^2\abs{A}^2 + \abs{\delta}^2\abs{B}^2 - \abs{\overline{\alpha}\gamma A + \overline{\beta}\delta B}^2 \\
        &= (1-\abs{\alpha}^2)\abs{\gamma}^2\abs{A}^2 + (1-\abs{\beta}^2)\abs{\delta}^2\abs{B}^2 - 2\Re\!\left(\overline{\alpha}\beta\gamma\overline{\delta}\,A\overline{B}\right) \\
        &= \abs{\beta}^2\abs{\gamma}^2\abs{A}^2 + \abs{\alpha}^2\abs{\delta}^2\abs{B}^2 - 2\Re\!\left(\overline{\alpha}\beta\gamma\overline{\delta}\,A\overline{B}\right) = \abs{\beta\gamma A - \alpha\delta B}^2.
\end{align*}
Again using $\bra{\psi_{target}^{x1}} = \bra{\psi_{target}^{x0}}D^\dagger$, we have
\[
    \beta\gamma A - \alpha\delta B = \bra{\psi_{target}^{x0}}\!\left(\beta\gamma\,\ket{\psi_{lab}^{x0}} - \alpha\delta\,D^\dagger\ket{\psi_{lab}^{x1}}\right),
\]
so by Cauchy--Schwarz and \eqref{eqn:rej-formula},
\[
    \Delta(\ket{\psi_{lab}^x}:\ket{\psi_{target}^x}) = \abs{\beta\gamma A - \alpha\delta B}^2 \leq \norm{\beta\gamma\,\ket{\psi_{lab}^{x0}} - \alpha\delta\,D^\dagger\ket{\psi_{lab}^{x1}}}^2 = \Pr[\textsc{SubTest}\text{ rejects}],
\]
which is \eqref{eqn:reject}.
\end{proof}

With Theorem~\ref{thm:subtest} in hand, the remaining ingredient is the following telescoping identity.

\begin{proposition}\label{prop:expectation}
    It holds that
    \[
        \mathop{\mathbb{E}}_{x}\!\left[\Delta(\ket{\psi_{lab}^x}:\ket{\psi_{target}^x})\right] = \frac{1}{n}\cdot\left(1 - \abs{\braket{\psi_{target}}{\psi_{lab}}}^2\right).
    \]
    Here, $x$ is sampled by performing Steps 1--2 of : first choose a uniformly random $k\in[n]$, then measure the first $k-1$ qubits of $\ket{\psi_{lab}}$ in the computational basis to obtain $x$.
\end{proposition}

\begin{proof}
    For $0\leq k\leq n$, define the quantity
    \[
        \Phi^k \coloneqq \mathop{\mathbb{E}}_{\substack{x,\\ |x|=k}}\!\left[\abs{\braket{\psi_{target}^{x}}{\psi_{lab}^{x}}}^2\right],
    \]
    where $x$ is sampled by measuring the first $k$ qubits of $\ket{\psi_{lab}}$ in the computational basis. Then,
    \begin{align*}
        \mathop{\mathbb{E}}_{x}\!\left[\Delta(\ket{\psi_{lab}^x}:\ket{\psi_{target}^x})\right]
        &= \mathop{\mathbb{E}}_{k\in[n]}\!\left[\mathop{\mathbb{E}}_{\substack{x,\\ |x|=k-1}}\mathop{\mathbb{E}}_{b}\!\left[\abs{\braket{\psi_{target}^{xb}}{\psi_{lab}^{xb}}}^2\right] - \mathop{\mathbb{E}}_{\substack{x,\\ |x|=k-1}}\!\left[\abs{\braket{\psi_{target}^x}{\psi_{lab}^x}}^2\right]\right]\\
        &= \mathop{\mathbb{E}}_{k\in[n]}\!\left[\mathop{\mathbb{E}}_{\substack{x,\\ |x|=k}}\!\left[\abs{\braket{\psi_{target}^{x}}{\psi_{lab}^{x}}}^2\right] - \mathop{\mathbb{E}}_{\substack{x,\\ |x|=k-1}}\!\left[\abs{\braket{\psi_{target}^x}{\psi_{lab}^x}}^2\right]\right]\\
        &= \mathop{\mathbb{E}}_{k\in[n]}\!\left[\Phi^k - \Phi^{k-1}\right]\\
        &= \frac{1}{n}\cdot(\Phi^n - \Phi^0) = \frac{1}{n}\cdot\left(1 - \abs{\braket{\psi_{target}}{\psi_{lab}}}^2\right),
    \end{align*}
    where the last line follows by a telescoping sum.
\end{proof}

Combining the two results, Theorem~\ref{thm:subtest} gives
\begin{align*}
    \Pr[\text{ rejects}]
        &= \mathop{\mathbb{E}}_{x}\Pr\!\left[\textsc{SubTest}_{\mathcal{T}_x}\!\left(\ket{\psi_{lab}^x}:\ket{\psi_{target}^x}\right) \text{ rejects}\right]\\
        &\geq \mathop{\mathbb{E}}_{x}\!\left[\Delta(\ket{\psi_{lab}^x}:\ket{\psi_{target}^x})\right]
        = \frac{1}{n}\cdot\left(1 - \abs{\braket{\psi_{target}}{\psi_{lab}}}^2\right),
\end{align*}
as claimed. As for the acceptance probability,
\begin{align*}
    \Pr[\text{ accepts}]
        &= \mathop{\mathbb{E}}_{x}\Pr\!\left[\textsc{SubTest}_{\mathcal{T}_x}\!\left(\ket{\psi_{lab}^x}:\ket{\psi_{target}^x}\right) \text{ accepts}\right]\\
        &\geq \mathop{\mathbb{E}}_{x}\!\left[\abs{\braket{\psi_{target}^x}{\psi_{lab}^x}}^2\right]
        \geq \abs{\braket{\psi_{target}}{\psi_{lab}}}^2,
\end{align*}
where the last inequality follows by Proposition~\ref{rem:nonneg}.

